Below are **only sections (2)–(11)**, written as constructive model definitions and numerical algorithms. No contracts, ledgers, certificates, or validation policy. This fills in the missing composition-to-generator part: species library → state basis → mobility/covariance → transitions → committor/Doob path moments → atmosphere/high-concentration/additive physics. It uses the same primitive heads already required by the conductivity-effect checklist: (c_i), (K_{ij}/Q_{ij}), (d_{ij}/M_{ij}), (D_i/R_i), and (A,h). 

---

```latex
\section{2. Complete species-library schema}

The electrolyte composition is
\[
x=(T,\{C_a\}_{a\in\mathcal C},\{\phi_s\}_{s\in\mathcal S},\mathcal L).
\]

Here \(T\) is temperature, \(C_a\) are conserved analytical component concentrations,
\(\phi_s\) are liquid fractions, and \(\mathcal L\) is the physical species library.
The library is
\[
\mathcal L=
\left(
\{\mathcal T_s\}_{s},
\{\mathcal Q_s\}_{s},
\{\mathcal G_s\}_{s},
\{\Theta_{st}\}_{s,t},
\mathcal M,
\mathcal B,
\mathcal R,
\mathcal \Psi
\right).
\]

For every species \(s\), the topology record is
\[
\mathcal T_s
=
(\mathcal I_s,\mathcal E_s^{\rm bond},\mathcal E_s^{\rm angle},
\mathcal E_s^{\rm torsion},\mathcal C_s^{\rm constraint}),
\]
where
\[
\mathcal I_s=\{1,\ldots,n_s\}
\]
is the site set.

For each site \(u\in\mathcal I_s\), define
\[
m_{su}>0,
\quad
r_{su}^{\rm atom}>0,
\quad
a_{su}^{\rm hyd}>0,
\quad
v_{su}>0,
\quad
\sigma_{su}>0,
\quad
\epsilon_{su}\ge0,
\quad
z_{su}\in\mathbb R,
\quad
\alpha_{su}\ge0.
\]

Meanings:
\[
m_{su}=\text{mass},
\]
\[
r_{su}^{\rm atom}=\text{steric radius},
\]
\[
a_{su}^{\rm hyd}=\text{hydrodynamic radius},
\]
\[
v_{su}=\text{site volume},
\]
\[
\sigma_{su},\epsilon_{su}=\text{Lennard--Jones parameters},
\]
\[
z_{su}=\text{formal or partial charge number},
\]
\[
\alpha_{su}=\text{polarizability}.
\]

The molecular charge-site model is
\[
\mathcal Q_s=\{(u,z_{su},r_{su}^{\rm charge},a_{su}^{\rm cloud})\}_{u\in\mathcal I_s},
\]
where \(a_{su}^{\rm cloud}\) is the finite charge-cloud radius.

The molecular geometry-feature record is
\[
\mathcal G_s=
(
V_s,\,
A_s^{\rm SASA},\,
a_s^{\rm hyd},\,
a_s^{\rm Born},\,
a_s^{\rm cloud},\,
\ell_s^{\rm asym},\,
n_s^{\rm donor},\,
n_s^{\rm acceptor},\,
n_s^{\rm HBA},\,
n_s^{\rm HBD},\,
\{g_{s\ell}^{\rm coord}\}_{\ell}
).
\]

Here
\[
V_s=\text{molecular volume},
\]
\[
A_s^{\rm SASA}=\text{solvent-accessible surface area},
\]
\[
a_s^{\rm hyd}=\text{hydrodynamic radius},
\]
\[
a_s^{\rm Born}=\text{Born cavity radius},
\]
\[
a_s^{\rm cloud}=\text{charge-cloud radius},
\]
\[
\ell_s^{\rm asym}=\text{ligand-field/shape asymmetry},
\]
\[
n_s^{\rm donor},n_s^{\rm acceptor},n_s^{\rm HBA},n_s^{\rm HBD}
=\text{coordination and H-bond descriptors},
\]
\[
g_{s\ell}^{\rm coord}
=\text{site-specific coordination free energy for ligand class }\ell.
\]

For every species pair \((s,t)\), define the interaction record
\[
\Theta_{st}
=
(
\Theta_{st}^{\rm LJ},
\Theta_{st}^{\rm Coul},
\Theta_{st}^{\rm Born},
\Theta_{st}^{\rm coord},
\Theta_{st}^{\rm PMF},
\Theta_{st}^{\rm mobility}
).
\]

Short-range mixing:
\[
\sigma_{uv}^{st}
=
\frac{\sigma_{su}+\sigma_{tv}}{2},
\qquad
\epsilon_{uv}^{st}
=
\sqrt{\epsilon_{su}\epsilon_{tv}}.
\]

Coulomb convention:
\[
U_{\rm Coul}
=
N_A
\frac{1}{4\pi\epsilon_0\epsilon_\infty}
\left(
U_{\rm real}+U_{\rm reciprocal}+U_{\rm self}+U_{\rm surface}
\right),
\]
with Ewald parameter \(\alpha_{\rm Ewald}\), real-space cutoff \(r_{\rm real}\),
reciprocal cutoff \(k_{\max}\), and boundary convention fixed by \(\Theta_{st}^{\rm Coul}\).

Born/desolvation term for bringing charge site \(u\) into environment \(e\):
\[
\Delta G_{su}^{\rm Born}(e)
=
-\frac{N_A e^2 z_{su}^2}{8\pi\epsilon_0 a_{su}^{\rm Born}}
\left(
1-\frac{1}{\epsilon_e}
\right).
\]

Pair desolvation penalty:
\[
\Delta G_{st}^{\rm desolv}
=
\sum_{u\in s}\sum_{v\in t}
\omega_{uv}^{\rm contact}
\left[
\Delta G_{su}^{\rm Born}(e_{\rm paired})
+
\Delta G_{tv}^{\rm Born}(e_{\rm paired})
-
\Delta G_{su}^{\rm Born}(e_{\rm bulk})
-
\Delta G_{tv}^{\rm Born}(e_{\rm bulk})
\right].
\]

Coordination energy:
\[
U_{\rm coord}(q)
=
N_A
\sum_{a,L}
\Delta g_{aL}^{\rm coord}\,
n_{aL}(q),
\]
where
\[
n_{aL}(q)=\sum_{b\in L}s_{aL}(r_{ab}),
\]
\[
s_{aL}(r)=\frac{1}{1+(r/r_{aL}^{0})^{p_{aL}}}.
\]

The mobility library \(\mathcal M\) defines a resistance contribution for each physical source:
\[
R_x(q)
=
R_{\rm Stokes}
+
R_{\rm hydro}
+
R_{\rm shape}
+
R_{\rm free\ volume}
+
R_{\rm charge\ cloud}
+
R_{\rm atmosphere}
+
R_{\rm cage}
+
R_{\rm constraint}.
\]

The diffusion tensor is
\[
D_x(q)=k_BT\,R_x(q)^\#.
\]

The basin library \(\mathcal B\) defines coordinate functions and thresholds for
free, SSIP, additive-separated SSIP, CIP, aggregate, ligand-shell, bridge,
orientation, cage, atmosphere, and local-environment basins.

The transition library \(\mathcal R\) defines allowed transition coordinates
\[
\xi_{ij}(q),
\]
transition domains
\[
\Omega_{ij},
\]
transition boundaries
\[
\partial A_i,\partial A_j,
\]
and committor boundary conditions.

The memory library \(\mathcal \Psi\) defines continuous slow variables
\[
\psi_\alpha(q)
\]
for atmosphere, cage, partner residence, ligand residence, anion orientation,
charge density, and local free-volume/stress relaxation.
```

---

```latex
\section{3. Deterministic basis-generation algorithm}

Define a feature map
\[
\mathcal F(q)=
\left(
f_{\rm pair}(q),
f_{\rm LiShell}(q),
f_{\rm ligand}(q),
f_{\rm anion}(q),
f_{\rm orient}(q),
f_{\rm cluster}(q),
f_{\rm env}(q),
f_{\rm atm}(q)
\right).
\]

Each component is computed from physical coordinates, not species names.

\subsection{Pair basin}

For Li center \(\ell\) and anion center \(a\),
\[
r_{\ell a}(q)=|r_\ell-r_a|.
\]

Define thresholds
\[
r_{\rm CIP}<r_{\rm SSIP}<r_{\rm free}.
\]

Then
\[
f_{\rm pair}(q)=
\begin{cases}
{\rm CIP}, & r_{\ell a}<r_{\rm CIP},\\
{\rm SSIP}, & r_{\rm CIP}\le r_{\ell a}<r_{\rm SSIP},\\
{\rm free}, & r_{\ell a}\ge r_{\rm free},\\
{\rm additive\_separated\_SSIP}, &
r_{\rm CIP}\le r_{\ell a}<r_{\rm SSIP}
\text{ and } n_{\ell L}\ge n_L^{\rm cut}.
\end{cases}
\]

\subsection{Li solvation shell}

For ligand class \(L\),
\[
n_{\ell L}(q)=\sum_{b\in L}s_{\ell L}(r_{\ell b}).
\]

Let
\[
n_{\ell{\rm solv}}=\sum_{L\in{\rm solvent}}n_{\ell L},
\qquad
n_{\ell{\rm add}}=\sum_{L\in{\rm neutral\ additive}}n_{\ell L},
\qquad
n_{\ell{\rm anion}}=\sum_{a\in{\rm anion}}s_{\ell a}(r_{\ell a}).
\]

Then
\[
f_{\rm LiShell}(q)=
\begin{cases}
{\rm solvent\_only}, & n_{\ell{\rm add}}<n_{\rm add}^{\rm cut}
\text{ and }n_{\ell{\rm anion}}<n_{\rm anion}^{\rm cut},\\
{\rm neutral\_ligand\_bound}, & n_{\ell{\rm add}}\ge n_{\rm add}^{\rm cut},\\
{\rm anion\_coordinated}, & n_{\ell{\rm anion}}\ge n_{\rm anion}^{\rm cut},\\
{\rm overcrowded}, & n_{\ell{\rm solv}}+n_{\ell{\rm add}}+n_{\ell{\rm anion}}
>n_{\rm shell}^{\max}.
\end{cases}
\]

\subsection{Neutral ligand state}

Let \(L\) be a neutral coordinating molecule class. Then
\[
f_{\rm ligand}(q)=
\begin{cases}
{\rm none}, & n_{\ell L}<n_L^{\rm cut},\\
{\rm monodentate\_bound}, & n_{\ell L}\ge n_L^{\rm cut}
\text{ and }m_L^{\rm dent}=1,\\
{\rm multidentate\_bound}, & n_{\ell L}\ge n_L^{\rm cut}
\text{ and }m_L^{\rm dent}>1,\\
{\rm ligand\_bridge}, &
n_{\ell_1 L}\ge n_L^{\rm cut}
\text{ and }n_{\ell_2 L}\ge n_L^{\rm cut}.
\end{cases}
\]

\subsection{Anion feature bin}

For anion \(a\), define descriptors
\[
a_a^{\rm hyd},\quad a_a^{\rm cloud},\quad \ell_a^{\rm asym},\quad m_a^{\rm dent}.
\]

Then
\[
f_{\rm anion}(a)=
\left(
b_{\rm hyd}(a_a^{\rm hyd}),
b_{\rm cloud}(a_a^{\rm cloud}),
b_{\rm asym}(\ell_a^{\rm asym}),
b_{\rm dent}(m_a^{\rm dent})
\right).
\]

Example binning:
\[
b_{\rm hyd}(a)=
\begin{cases}
{\rm compact}, & a<a_{\rm hyd}^{(1)},\\
{\rm medium}, & a_{\rm hyd}^{(1)}\le a<a_{\rm hyd}^{(2)},\\
{\rm bulky}, & a\ge a_{\rm hyd}^{(2)}.
\end{cases}
\]

\[
b_{\rm cloud}(a)=
\begin{cases}
{\rm localized}, & a<a_{\rm cloud}^{(1)},\\
{\rm delocalized}, & a\ge a_{\rm cloud}^{(1)}.
\end{cases}
\]

\[
b_{\rm dent}(m)=
\begin{cases}
{\rm monodentate}, & m=1,\\
{\rm multidentate}, & m>1.
\end{cases}
\]

\subsection{Anion orientation}

Let \(u_a(q)\) be the principal axis of the anion and
\[
\hat r_{\ell a}=\frac{r_a-r_\ell}{|r_a-r_\ell|}.
\]
Define
\[
\omega_a(q)=u_a(q)\cdot \hat r_{\ell a}.
\]

Then
\[
f_{\rm orient}(q)=
\begin{cases}
{\rm radial}, & \omega_a>\omega_+,\\
{\rm tangential}, & |\omega_a|\le\omega_0,\\
{\rm bridging}, & \omega_a<-\omega_-,\\
{\rm free\_rotating}, & \tau_{\rm orient}<\tau_{\rm orient}^{\rm cut}.
\end{cases}
\]

\subsection{Cluster state}

Construct a graph
\[
G(q)=(V,E(q)),
\]
where vertices are charged centers and
\[
(a,b)\in E(q)
\iff
r_{ab}<r_{\rm pair}^{\rm cut}.
\]

Connected components define clusters. For cluster \(C\):
\[
N_{\rm Li}(C)=\sum_{v\in C}\mathbf 1_{\{v={\rm Li}\}},
\]
\[
N_A(C)=\sum_{v\in C}\mathbf 1_{\{v={\rm anion}\}},
\]
\[
z_C=\sum_{v\in C}z_v.
\]

Then
\[
f_{\rm cluster}(C)=
\left(N_{\rm Li}(C),N_A(C),z_C,\tau_C,\text{topology}(C)\right).
\]

Cluster labels include
\[
{\rm Li^+},\quad
{\rm A^-},\quad
{\rm LiA},\quad
{\rm Li_2A^+},\quad
{\rm LiA_2^-},\quad
{\rm Li_2A_2},\quad
{\rm higher\_cluster},\quad
{\rm bridge\_network}.
\]

\subsection{Local environment}

Define local packing
\[
\varphi(q)=\sum_b v_b W(|r_b-r_0|),
\]
local ionic strength
\[
I(q)=\frac12\sum_i z_i^2 c_i^{\rm loc}(q),
\]
local dielectric
\[
\epsilon_{\rm loc}(q)=\sum_s \phi_s^{\rm loc}(q)\epsilon_s,
\]
and local viscosity
\[
\eta_{\rm loc}(q)=\eta_{\rm bulk}\,
f_{\rm salt}(I(q))f_{\rm free\ volume}(\varphi(q))f_{\rm additive}(q).
\]

Then
\[
f_{\rm env}(q)=
\left(
b_\varphi(\varphi),
b_I(I),
b_\epsilon(\epsilon_{\rm loc}),
b_\eta(\eta_{\rm loc})
\right).
\]

\subsection{Joint sparse state}

The full projected state key is
\[
\alpha(q)=
\left[
f_{\rm pair},
f_{\rm LiShell},
f_{\rm ligand},
f_{\rm anion},
f_{\rm orient},
f_{\rm cluster},
f_{\rm env},
f_{\rm atm}
\right].
\]

The basin is
\[
A_\alpha=\{q:\alpha(q)=\alpha\}.
\]

Instantiate \(A_\alpha\) only if
\[
Z_\alpha=
\int_{A_\alpha}e^{-U(q)/(RT)}\,dq>Z_{\min}
\]
and the stoichiometry is feasible:
\[
\nu_{\alpha,a}\le N_a
\qquad
\forall a.
\]

Thus the basis is sparse:
\[
\mathcal A=\{A_\alpha:Z_\alpha>Z_{\min}\}.
\]
```

This directly instantiates the states the checklist says are missing: SSIP/CIP/aggregate populations, ligand-conditioned Li solvation shells, additive-bound Li states, anion coordination/orientation states, and high-viscosity/low-dielectric regimes. The source material explicitly requires shared Li, SSIP/CIP, Li-additive, aggregate motifs, and name-blind site-measure state generation. 

---

```latex
\section{4. Physical state classification and DC transport projector}

For each basin \(A_\alpha\), define the charged centers
\[
\mathcal C_\alpha=\{p=1,\ldots,m_\alpha\}.
\]

Each center has
\[
z_{\alpha p}\in\mathbb R,
\quad
r_{\alpha p}\in\mathbb R^3,
\quad
a_{\alpha p}^{\rm hyd},
\quad
a_{\alpha p}^{\rm cloud},
\quad
S_{\alpha p}^{\rm shape}.
\]

Let
\[
z_\alpha=(z_{\alpha1},\ldots,z_{\alpha m_\alpha})^\top.
\]

Define the coordinate-space projector
\[
\Pi_\alpha(q)\in\mathbb R^{d\times d}
\]
that selects unbounded DC charge-transport directions.

\subsection{Free ion}

For free ion \(p\),
\[
\Pi_\alpha=\Pi_{\rm trans}(p),
\]
where \(\Pi_{\rm trans}(p)\) projects onto the three unwrapped translational coordinates of \(p\).

\subsection{Charged cluster COM}

For a charged cluster with net charge
\[
z_{\rm net,\alpha}=\sum_p z_{\alpha p}\ne0,
\]
the cluster COM coordinate is
\[
R_{\rm COM,\alpha}=
\frac{\sum_p m_{\alpha p}r_{\alpha p}}
{\sum_p m_{\alpha p}}.
\]

The COM projector is
\[
\Pi_\alpha=\Pi_{\rm COM,\alpha}
\]
only if the cluster translation is unbounded.

\subsection{Neutral cluster COM}

If
\[
z_{\rm net,\alpha}=0
\]
and all charged centers move as a rigid co-moving neutral cluster, then
\[
G_P\Pi_{\rm COM,\alpha}=0.
\]
Thus the neutral COM contributes no DC self-current.

\subsection{Bound internal polarization}

For a bounded internal coordinate \(\xi_{\rm int}\),
\[
\Pi_\alpha\xi_{\rm int}=0.
\]
Bound internal polarization contributes to \(A,h\), not \(D_\alpha^{\rm self}\).

\subsection{Associated charge-identity diffusion}

If a bound/associated state permits unbounded charge identity transfer,
define a charge-identity coordinate \(y_a\). Then
\[
\Pi_\alpha=\Pi_{\rm id,a}
\]
only along \(y_a\).

The associated identity diffusivity is
\[
D_{{\rm id},a}
=
\frac{1}{6}
k_{{\rm switch},a}
p_{{\rm switch},a}
\mathbb E[\ell_{{\rm switch},a}^{2}].
\]

Here
\[
k_{{\rm switch},a}
\]
is the validated switching rate,
\[
p_{{\rm switch},a}
\]
is the probability that the charge identity changes partner rather than returning,
and
\[
\ell_{{\rm switch},a}
\]
is the unwrapped charge-identity displacement.

If \(k_{\rm switch}\), \(p_{\rm switch}\), or \(\ell_{\rm switch}\) is unavailable,
\[
D_{{\rm id},a}=0.
\]

\subsection{State self-current}

The within-state tensor is
\[
D_\alpha^{\rm self}
=
\frac{1}{c_\alpha}
\int_{A_\alpha}
G_P(q)\Pi_\alpha(q)D_x(q)\Pi_\alpha(q)^\top G_P(q)^\top
\rho(dq).
\]

For center covariance representation, define
\[
D_\alpha^{\rm centers}\in\mathbb R^{m_\alpha\times m_\alpha}.
\]

Then
\[
D_Q(\alpha)=z_\alpha^\top D_\alpha^{\rm centers}z_\alpha.
\]

For Li--anion:
\[
D_Q
=
D_{\rm Li}
+
D_{\rm A}
-
2D_{\rm Li,A}.
\]

If
\[
D_{\rm Li,A}<0,
\]
then
\[
-2D_{\rm Li,A}>0,
\]
so anti-correlated Li--anion motion increases conductivity.
```

This is the non-negotiable fix for the recurring neutral-pair/SSIP issue: the model must use state-resolved charged-center tensors and not collapse multi-charge motifs to scalar net charge. That requirement appears explicitly in the finite-generator plan and the effect checklist.  

---

```latex
\section{5. Complete transition-family enumeration}

Every transition is an ordered edge
\[
e=(i,j,\kappa)
\]
from basin \(A_i\) to basin \(A_j\), with transition family \(\kappa\).

Each edge has
\[
K_{ij}^{(\kappa)}=K_{ji}^{(\kappa)}\ge0,
\]
\[
Q_{ij}^{(\kappa)}=\frac{K_{ij}^{(\kappa)}}{c_i},
\]
\[
d_{ji}^{(\kappa)}=-d_{ij}^{(\kappa)},
\]
\[
M_{ji}^{(\kappa)}=M_{ij}^{(\kappa)}.
\]

The total flux is
\[
K_{ij}=\sum_\kappa K_{ij}^{(\kappa)}.
\]

The total generator is
\[
Q_{ij}=\frac{K_{ij}}{c_i}.
\]

\subsection{Free \(\leftrightarrow\) SSIP}

Reaction coordinate:
\[
\xi_{\rm free,SSIP}=r_{\rm LiA}.
\]

Basins:
\[
A_{\rm free}=\{r_{\rm LiA}\ge r_{\rm SSIP}\},
\]
\[
A_{\rm SSIP}=\{r_{\rm CIP}\le r_{\rm LiA}<r_{\rm SSIP}\}.
\]

Transition moment:
\[
\Delta P=P(q_{\rm SSIP})-P(q_{\rm free}).
\]

If the transition only changes association label without unwrapped charge motion:
\[
d=0,\qquad M=0.
\]

If the transition includes a structural approach/separation hop:
\[
d=\mathbb E[\Delta P],
\qquad
M=\mathbb E[\Delta P\Delta P^\top].
\]

\subsection{SSIP \(\leftrightarrow\) CIP}

Reaction coordinate:
\[
\xi_{\rm SSIP,CIP}=r_{\rm LiA}.
\]

Basins:
\[
A_{\rm CIP}=\{r_{\rm LiA}<r_{\rm CIP}\}.
\]

Pure conversion:
\[
d=0,\qquad M=0.
\]

Charge-carrying structural conversion:
\[
d,M
\]
must be computed from transition-path \(\Delta P\) samples.

\subsection{Free/SSIP/CIP \(\leftrightarrow\) aggregate}

Cluster graph transition:
\[
G(q)\to G'(q)
\]
with changed connected component size.

Coordinates:
\[
\xi_{\rm agg}=
\left(
r_{\rm LiA_1},\ldots,r_{\rm LiA_k},
N_{\rm cluster},
z_{\rm cluster}
\right).
\]

Edges:
\[
{\rm free}\leftrightarrow{\rm Li_2A^+},
\quad
{\rm free}\leftrightarrow{\rm LiA_2^-},
\quad
{\rm SSIP}\leftrightarrow{\rm Li_2A_2},
\quad
{\rm CIP}\leftrightarrow{\rm bridge}.
\]

Pure cluster relabeling:
\[
d=0,\quad M=0.
\]

Cluster-mediated charge transfer:
\[
d,M
\]
from physical charge-displacement samples only.

\subsection{Li \(\leftrightarrow\) Li--ligand}

Reaction coordinate:
\[
\xi_{\rm Li,L}=n_{\rm LiL}.
\]

Basins:
\[
A_{\rm Li}=\{n_{\rm LiL}<n_L^{\rm cut}\},
\]
\[
A_{\rm LiL}=\{n_{\rm LiL}\ge n_L^{\rm cut}\}.
\]

Pure ligand binding:
\[
d=0,\quad M=0.
\]

If ligand exchange produces unbounded Li shell hopping:
\[
d,M
\]
from \(\Delta P\) samples.

\subsection{Li--anion \(\leftrightarrow\) Li--ligand--anion}

Reaction coordinate:
\[
\xi_{\rm ternary}=(r_{\rm LiA},n_{\rm LiL}).
\]

Basins:
\[
A_{\rm LiA}
=
\{r_{\rm LiA}<r_{\rm SSIP},\,n_{\rm LiL}<n_L^{\rm cut}\},
\]
\[
A_{\rm LiL\cdots A}
=
\{r_{\rm CIP}\le r_{\rm LiA}<r_{\rm SSIP},\,n_{\rm LiL}\ge n_L^{\rm cut}\}.
\]

This transition changes:
\[
c_i,\quad K_{ij},\quad D_{\rm Li,A|i},\quad A,h.
\]

\subsection{Ordinary SSIP \(\leftrightarrow\) additive-separated SSIP}

Reaction coordinate:
\[
\xi_{\rm addSSIP}=(r_{\rm LiA},n_{\rm LiL},\chi_L),
\]
where \(\chi_L\) identifies ligand insertion between Li and anion.

Pure insertion:
\[
d=0,\quad M=0.
\]

If insertion triggers partner switch:
\[
d,M
\]
from \(\Delta P\).

\subsection{Partner switch}

Old partner \(a\), new partner \(b\):
\[
{\rm Li}\cdots a \to {\rm Li}\cdots b.
\]

Reaction coordinate:
\[
\xi_{\rm switch}=r_{\rm Li,a}-r_{\rm Li,b}.
\]

Charge displacement:
\[
\Delta P=P(q_{\rm after})-P(q_{\rm before}).
\]

This edge always requires:
\[
K>0,\quad d\ne0\text{ or }M\ne0
\]
only if the transition changes unwrapped charge identity.

\subsection{Identity diffusion}

For charge identity coordinate \(y\):
\[
y\to y+\ell.
\]

Moments:
\[
d=\mathbb E[\ell],
\qquad
M=\mathbb E[\ell\ell^\top].
\]

\subsection{Structural hop}

A structural hop is legal only when a physical displacement coordinate exists:
\[
\Delta P=P(q_{\rm end})-P(q_{\rm start}).
\]

Then:
\[
d=\mathbb E[\Delta P],
\qquad
M=\mathbb E[\Delta P\Delta P^\top].
\]

\subsection{Cage capture/release}

Cage state:
\[
A_{\rm cage}
=
\{|\psi_{\rm cage}|>\psi_{\rm cage}^{\rm cut}\}.
\]

Capture/release without charge displacement:
\[
d=0,\quad M=0.
\]

Backjump memory enters through either finite-state \(C^Q\) or Mori \(A,h\).

\subsection{Atmosphere capture/release}

If atmosphere is represented discretely:
\[
A_{\rm atm,k}=\{p_{\rm atm}\in[p_k,p_{k+1})\}.
\]

Pure atmosphere relaxation:
\[
d=0,\quad M=0.
\]

Current coupling enters through Mori memory:
\[
A_{\alpha\gamma},\quad h_{\alpha a}.
\]
```

---

```latex
\section{6. Committor PDE numerical policy}

For each edge \(i\leftrightarrow j\), define transition domain
\[
\Omega_{ij}=\Omega\setminus(A_i\cup A_j).
\]

The committor solves
\[
Lq_{ij}=0
\quad\text{in }\Omega_{ij},
\]
\[
q_{ij}=0\quad\text{on }\partial A_i,
\qquad
q_{ij}=1\quad\text{on }\partial A_j.
\]

\subsection{Finite-element discretization}

Construct a mesh
\[
\mathcal T_{ij}^{(r)}
\]
on \(\Omega_{ij}\), where \(r\) is refinement level.

Use continuous piecewise-linear basis functions
\[
\{\varphi_m\}_{m=1}^{M_r}.
\]

Write
\[
q_{ij}^{(r)}(q)
=
q_D(q)+\sum_{m=1}^{M_r}u_m\varphi_m(q),
\]
where \(q_D\) enforces Dirichlet values.

Stiffness matrix:
\[
S_{mn}^{(r)}
=
\int_{\Omega_{ij}}
\nabla\varphi_m^\top D(q)\nabla\varphi_n\,\rho(dq).
\]

Right-hand side:
\[
f_m^{(r)}
=
-
\int_{\Omega_{ij}}
\nabla\varphi_m^\top D(q)\nabla q_D\,\rho(dq).
\]

Solve:
\[
S^{(r)}u^{(r)}=f^{(r)}.
\]

Use conjugate gradient if \(S^{(r)}\) is symmetric positive definite. Use
MINRES if constraints produce semidefinite blocks. Precondition with incomplete
Cholesky or algebraic multigrid.

\subsection{Capacity computation}

Compute
\[
K_{ij}^{(r)}
=
\int_{\Omega_{ij}}
\nabla q_{ij}^{(r)\top}D(q)\nabla q_{ij}^{(r)}\,\rho(dq).
\]

\subsection{Adaptive refinement}

Refine elements with largest local energy residual
\[
\eta_E
=
\int_E
\nabla q_{ij}^{(r)\top}D\nabla q_{ij}^{(r)}\,\rho(dq).
\]

Refine the smallest set of elements \(\mathcal R\) satisfying
\[
\sum_{E\in\mathcal R}\eta_E
\ge
\theta_{\rm refine}\sum_E\eta_E,
\]
with
\[
0.3\le\theta_{\rm refine}\le0.7.
\]

Re-solve until
\[
\frac{|K_{ij}^{(r+1)}-K_{ij}^{(r)}|}
{\max(K_{ij}^{(r)},K_{\min})}
<
\tau_K.
\]

Default:
\[
\tau_K=10^{-3},
\qquad
K_{\min}=10^{-30}\ \mathrm{mol\,m^{-3}\,s^{-1}}.
\]

\subsection{Disconnected domains}

If \(\Omega_{ij}\) has disconnected components \(\Omega_{ij}^{(m)}\), solve
on each component. If a component touches both boundaries, include it. If it
touches only one boundary, it contributes zero flux.

Thus:
\[
K_{ij}=\sum_mK_{ij}^{(m)}.
\]

\subsection{Tiny populations}

If
\[
c_i<c_{\min}
\]
or
\[
c_j<c_{\min},
\]
remove the edge from the active generator and merge the basin or mark it inactive.
Default:
\[
c_{\min}=10^{-18}\ \mathrm{mol\,m^{-3}}.
\]
```

---

```latex
\section{7. Doob-conditioned transition-path moment solver}

For each edge \(i\to j\), use the converged committor \(q_{ij}\).

Original overdamped drift:
\[
b(q)=
\nabla\cdot D(q)
-
\frac{1}{RT}D(q)\nabla U(q).
\]

Doob-conditioned drift:
\[
b^{(ij)}(q)
=
b(q)
+
2D(q)\nabla\log q_{ij}(q).
\]

Regularize:
\[
\nabla\log q_{ij}
=
\frac{\nabla q_{ij}}{\max(q_{ij},q_{\min})}.
\]

Default:
\[
q_{\min}=10^{-8}.
\]

\subsection{Exit-surface sampling}

Exit density:
\[
p_{ij}^{\rm exit}(q)
=
\frac{
\rho(q)\,n_{ij}(q)^\top D(q)\nabla q_{ij}(q)
}{
K_{ij}
}.
\]

Sample exit points by weighted surface quadrature:
\[
\Pr(q=q_\ell^\Sigma)
=
\frac{
w_\ell^\Sigma\rho(q_\ell^\Sigma)n^\top D(q_\ell^\Sigma)\nabla q_{ij}(q_\ell^\Sigma)
}{
\sum_m
w_m^\Sigma\rho(q_m^\Sigma)n^\top D(q_m^\Sigma)\nabla q_{ij}(q_m^\Sigma)
}.
\]

\subsection{Time-step rule}

At point \(q\), choose
\[
\Delta t(q)
=
\min\left[
\Delta t_{\max},
\gamma_t
\frac{\ell_{\rm mesh}(q)^2}{\lambda_{\max}(D(q))},
\gamma_b
\frac{\ell_{\rm mesh}(q)}{\|b^{(ij)}(q)\|}
\right].
\]

Defaults:
\[
\gamma_t=10^{-2},
\qquad
\gamma_b=10^{-1}.
\]

\subsection{Euler--Maruyama propagation}

At step \(n\):
\[
q_{n+1}
=
q_n
+
b^{(ij)}(q_n)\Delta t_n
+
L_D(q_n)\sqrt{2\Delta t_n}\,\eta_n,
\]
where
\[
D(q_n)=L_D(q_n)L_D(q_n)^\top,
\]
and
\[
\eta_n\sim\mathcal N(0,I).
\]

Stop at first hit:
\[
\tau_j=\inf\{t:q_t\in A_j\}.
\]

If a step crosses the boundary, linearly interpolate:
\[
q_{\rm hit}=q_n+\theta(q_{n+1}-q_n),
\]
where \(\theta\in[0,1]\) solves the boundary equation.

\subsection{Path charge displacement}

For each path:
\[
\Delta P_\ell=P(q_{\rm hit})-P(q_0).
\]

Reject path if
\[
\|\Delta P_\ell\|>L_{\max}.
\]

Default:
\[
L_{\max}=10^{-8}\ \mathrm m.
\]

\subsection{Moment estimates}

After \(N\) accepted paths:
\[
d_{ij}^{(N)}
=
\frac{1}{N}\sum_{\ell=1}^N\Delta P_\ell,
\]
\[
M_{ij}^{(N)}
=
\frac{1}{N}\sum_{\ell=1}^N\Delta P_\ell\Delta P_\ell^\top.
\]

Block convergence: split paths into \(B\) blocks. Compute block means
\[
d_{ij}^{[b]},\quad M_{ij}^{[b]}.
\]

Stop when
\[
\frac{\operatorname{SE}(d_{ij})}{\max(\|d_{ij}\|,d_{\min})}<\tau_d
\]
and
\[
\frac{\operatorname{SE}(M_{ij})}{\max(\|M_{ij}\|_F,M_{\min})}<\tau_M.
\]

Defaults:
\[
\tau_d=0.05,
\qquad
\tau_M=0.05.
\]
```

---

```latex
\section{8. State-local mobility and charged-center covariance construction}

For each state \(\alpha\), define centers
\[
p=1,\ldots,m_\alpha.
\]

Each center has charge \(z_{\alpha p}\), position \(r_{\alpha p}\), hydrodynamic
radius \(a_{\alpha p}^{\rm hyd}\), charge-cloud radius \(a_{\alpha p}^{\rm cloud}\),
shape tensor \(S_{\alpha p}\), and relative position
\[
\rho_{\alpha p}=r_{\alpha p}-R_\alpha^{\rm ref}.
\]

Construct total resistance:
\[
R_\alpha
=
R_\alpha^{\rm Stokes}
+
R_\alpha^{\rm hydro}
+
R_\alpha^{\rm shape}
+
R_\alpha^{\rm free\ volume}
+
R_\alpha^{\rm charge\ cloud}
+
R_\alpha^{\rm atmosphere}
+
R_\alpha^{\rm constraint}
+
R_\alpha^{\rm cage}.
\]

\subsection{Stokes resistance}

For center \(p\):
\[
\zeta_{\alpha p}=6\pi\eta_\alpha a_{\alpha p}^{\rm hyd}.
\]

Then
\[
R_{\alpha,pp}^{\rm Stokes}=\zeta_{\alpha p}I_3.
\]

\subsection{Shape resistance}

Let \(S_{\alpha p}\) have eigenvalues \(s_1,s_2,s_3\). Define
\[
R_{\alpha p}^{\rm shape}
=
\zeta_{\alpha p}
\left[
I_3+\lambda_{\rm shape}(S_{\alpha p}-I_3)
\right].
\]

For scalar shape factor:
\[
R_{\alpha p}^{\rm shape}
=
\zeta_{\alpha p}\,
\chi_{\alpha p}^{\rm shape} I_3,
\]
\[
\chi_{\alpha p}^{\rm shape}
=
\left(1+\ell_{\alpha p}^{\rm asym}\right)^{\gamma_{\rm shape}}.
\]

\subsection{Free-volume resistance}

Local packing:
\[
\varphi_\alpha=\mathbb E[\varphi(q)\mid A_\alpha].
\]

Free-volume factor:
\[
\chi_\alpha^{\rm fv}
=
\left(1-\frac{\varphi_\alpha}{\varphi_{\max}}\right)^{-\gamma_{\rm fv}}.
\]

Then
\[
R_\alpha^{\rm free\ volume}
=
(\chi_\alpha^{\rm fv}-1)R_\alpha^{\rm Stokes}.
\]

\subsection{Charge-cloud resistance}

For center \(p\):
\[
\rho_{\alpha p}^{\rm charge}
=
\frac{|z_{\alpha p}|}{(a_{\alpha p}^{\rm cloud})^3}.
\]

Charge-cloud resistance:
\[
R_{\alpha p}^{\rm cloud}
=
\lambda_{\rm cloud}
\left(\rho_{\alpha p}^{\rm charge}\right)^{\gamma_{\rm cloud}}
F_{\rm screen}(\kappa a_{\alpha p}^{\rm cloud})
I_3.
\]

Screening form factor:
\[
F_{\rm screen}(x)=\frac{x^2}{1+x^2}.
\]

\subsection{Constraint resistance}

For constraint mode \(m\):
\[
R_{\alpha}^{\rm constraint,m}
=
k_BT
\frac{\tau_{\alpha m}}{\ell_{\alpha m}^2}
s_{\alpha m}s_{\alpha m}^\top.
\]

Here \(s_{\alpha m}\) is the relative-motion direction, \(\tau_{\alpha m}\)
is the residence time, and \(\ell_{\alpha m}\) is the constraint length.

\[
R_{\alpha}^{\rm constraint}=\sum_mR_{\alpha}^{\rm constraint,m}.
\]

\subsection{Cage resistance}

For cage coordinate direction \(s_{\rm cage}\):
\[
R_{\alpha}^{\rm cage}
=
k_BT
\frac{\tau_{\rm cage,\alpha}}{\ell_{\rm cage,\alpha}^2}
s_{\rm cage}s_{\rm cage}^\top.
\]

\subsection{Mobility and charge covariance}

The mobility is
\[
D_\alpha=k_BT\,R_\alpha^\#.
\]

Let
\[
Z_\alpha
\]
be the matrix mapping center velocities to charge displacement:
\[
\dot P=Z_\alpha v.
\]

Then
\[
D_\alpha^{Q}
=
Z_\alpha D_\alpha Z_\alpha^\top.
\]

For one Cartesian axis:
\[
D_Q(\alpha)=z_\alpha^\top D_\alpha^{\rm centers}z_\alpha.
\]

This is the state-local self-current tensor:
\[
D_\alpha^{\rm self}=D_\alpha^Q
\]
projected onto unbounded coordinates.
```

---

```latex
\section{9. Ion-atmosphere operator}

The atmosphere operator is a resistance matrix projected into state charge geometry.

\subsection{Screening}

Ionic strength:
\[
I=\frac12\sum_s z_s^2 c_s.
\]

Inverse Debye length:
\[
\kappa^2
=
\frac{F^2}{\epsilon_0\epsilon_r RT}
\sum_s z_s^2 c_s.
\]

Debye length:
\[
\lambda_D=\kappa^{-1}.
\]

\subsection{Bulk charge-mode response}

For wavevector \(k\), define charge-density form factor:
\[
\rho_z(k;q)=\sum_p z_p e^{ik\cdot r_p}.
\]

Gradient:
\[
g_k(q)=\nabla_q\rho_z(k;q).
\]

Define the bulk PNP/Stokes resistance response:
\[
\mathcal R_{\rm bulk}(k)
=
\mathcal R_{\rm ep}(k)+\mathcal R_{\rm rel}(k).
\]

Electrophoretic part:
\[
\mathcal R_{\rm ep}(k)
=
\Gamma_{\rm ep}
\frac{k^2}{(k^2+\kappa^2)^2}
\eta.
\]

Relaxation part:
\[
\mathcal R_{\rm rel}(k)
=
\Gamma_{\rm rel}
\frac{\kappa^2}{(k^2+\kappa^2)^2}
\frac{1}{D_{\rm amb}k^2+\tau_{\rm chem}^{-1}}.
\]

Ambipolar diffusivity:
\[
D_{\rm amb}
=
\frac{\sum_s z_s^2c_sD_s}{\sum_s z_s^2c_s}.
\]

\subsection{State charge-geometry form factor}

For state \(\alpha\), define
\[
F_\alpha(k)
=
\sum_{p\in\alpha} z_{\alpha p}
e^{ik\cdot \rho_{\alpha p}}
\exp\left[-\frac12k^2(a_{\alpha p}^{\rm cloud})^2\right].
\]

Then the state atmosphere resistance is
\[
R_\alpha^{\rm atmosphere}
=
\sum_{k\in\mathcal K}
|F_\alpha(k)|^2
\mathcal R_{\rm bulk}(k)
\,G_{\alpha k}G_{\alpha k}^\top.
\]

Here \(G_{\alpha k}\) maps the state center coordinate velocities into the
charge-density mode \(k\).

\subsection{Neutral-pair cancellation}

For a co-located neutral pair:
\[
z_1=+1,\quad z_2=-1,\quad \rho_1=\rho_2.
\]

Then
\[
F_\alpha(k)=e^{ik\cdot\rho_1}-e^{ik\cdot\rho_1}=0.
\]

Thus
\[
R_\alpha^{\rm atmosphere}=0
\]
for the neutral charge-neutral translation mode.

\subsection{Separated-pair interpolation}

For a Li--anion pair separated by \(r\):
\[
F_\alpha(k)=e^{ik\cdot r_{\rm Li}}-e^{ik\cdot r_{\rm A}}.
\]

At small \(kr\):
\[
|F_\alpha(k)|^2\approx (k\cdot r)^2.
\]

At large \(kr\), angular average gives
\[
\langle |F_\alpha(k)|^2\rangle\to2.
\]

Therefore:
\[
R_{\rm atmosphere}^{\rm SSIP/CIP}(r)
\]
interpolates between neutral cancellation and independent-ion resistance.

\subsection{Debye--Falkenhagen diagnostic timescale}

Atmosphere relaxation time:
\[
\tau_{\rm atm}=\frac{1}{D_{\rm amb}\kappa^2}.
\]

State lifetime:
\[
\tau_\alpha.
\]

Diagnostic gate:
\[
g_\alpha^{\rm diag}
=
\frac{\tau_\alpha}{\tau_\alpha+\tau_{\rm atm}}.
\]

This gate is diagnostic unless explicitly included in \(R_{\rm rel}\). If used:
\[
R_{\alpha}^{\rm rel,eff}
=
g_\alpha R_{\alpha}^{\rm rel}.
\]

It must not multiply final \(\sigma\).
```

This matches the uploaded plan: atmosphere resistance must be a finite-size nonideal carrier matrix response projected into state-local charge geometry, not a scalar trace multiplier; neutral pairs cancel, SSIP/contact pairs interpolate by Li–anion separation, and the Debye–Falkenhagen gate is diagnostic until geometry invariants pass. 

---

```latex
\section{10. High-concentration closure}

High concentration modifies \(U\), \(D\), \(K\), and \(A,h\). It does not modify
the final conductivity formula.

\subsection{Finite occupied volume}

Each species \(s\) has molecular volume \(V_s\). Local occupied volume fraction:
\[
\varphi(q)=\sum_s V_s c_s^{\rm loc}(q).
\]

Packing free energy:
\[
U_{\rm pack}(q)
=
RT
\int
\frac{\varphi(r)(4-3\varphi(r))}
{(1-\varphi(r))^2}
\,dr.
\]

If
\[
\varphi(r)\ge\varphi_{\max}
\]
for any quadrature point, the state has infinite free energy:
\[
U_{\rm pack}=+\infty.
\]

\subsection{Activity correction}

Debye--Huckel activity:
\[
\ln\gamma_i^{\rm DH}
=
-\frac{A_\gamma z_i^2\sqrt I}{1+B_\gamma a_i\sqrt I}.
\]

Hard-sphere activity:
\[
\ln\gamma_i^{\rm HS}
=
\frac{\partial}{\partial n_i}
\left[
\frac{\varphi(4-3\varphi)}{(1-\varphi)^2}
\right].
\]

Total activity contribution:
\[
U_{\rm act}
=
RT\sum_i n_i
\left(
\ln\gamma_i^{\rm DH}
+
\ln\gamma_i^{\rm HS}
\right).
\]

\subsection{Dielectric decrement}

Effective dielectric:
\[
\epsilon_r(c)
=
\epsilon_r^{0}
-
\sum_s a_s^{\epsilon}c_s
-
b_\epsilon I.
\]

Require
\[
\epsilon_r(c)\ge\epsilon_{\min}.
\]

This affects:
\[
U_{\rm Coul},\quad U_{\rm Born},\quad \kappa,\quad R_{\rm atmosphere}.
\]

\subsection{Jones--Dole viscosity}

Bulk viscosity:
\[
\eta(c)
=
\eta_0
\left[
1+
\sum_s A_s^{\eta}\sqrt{c_s}
+
\sum_s B_s^{\eta}c_s
+
\sum_{s,t}C_{st}^{\eta}c_sc_t
\right].
\]

Local microviscosity:
\[
\eta_{\rm loc}(q)
=
\eta(c)
\left(1-\frac{\varphi(q)}{\varphi_{\max}}\right)^{-\gamma_{\rm fv}}
\left[
1+
\sum_L a_L^{\eta}n_L(q)
\right].
\]

This enters:
\[
R_{\rm Stokes},\quad R_{\rm hydro},\quad K_{ij},\quad A.
\]

\subsection{Aggregate saturation}

For cluster \(C\) with stoichiometry \(\nu_C\), free energy:
\[
G_C
=
G_C^{0}
+
G_C^{\rm Coul}
+
G_C^{\rm desolv}
+
G_C^{\rm coord}
+
G_C^{\rm steric}
+
G_C^{\rm entropy}
+
G_C^{\rm activity}.
\]

Cluster entropy penalty:
\[
G_C^{\rm entropy}
=
RT\left[
\ln(n_C!)
+
\gamma_{\rm rot}\ln\Omega_C^{-1}
\right].
\]

Steric penalty:
\[
G_C^{\rm steric}
=
RT
\left(\frac{V_C}{V_{\rm void}}\right)^{\gamma_{\rm steric}}.
\]

Mass action:
\[
c_C
=
C^\circ
\exp[-G_C/(RT)]
\prod_a
\left(\frac{c_a}{C^\circ}\right)^{\nu_{C,a}}.
\]

At high concentration, aggregates reduce free-ion \(c_i\), alter \(D_i^{\rm self}\),
and create aggregate transition edges \(K_{ij}\).
```

---

```latex
\section{11. Additive and neutral-ligand handling}

Additives enter only as neutral ligand feature classes after species parsing.

For neutral ligand \(L\), define feature vector
\[
\ell_L=
(
n_L^{\rm donor},
n_L^{\rm acceptor},
n_L^{\rm HBA},
n_L^{\rm HBD},
\alpha_L,
V_L,
a_L^{\rm hyd},
g_{\rm LiL}^{\rm coord},
\chi_L^{\rm visc},
m_L^{\rm dent}
).
\]

No equation below depends on the string name of the molecule.

\subsection{Li--ligand coordination}

Coordination number:
\[
n_{\rm Li,L}(q)=\sum_{b\in L}s_{\rm Li,L}(r_{\rm Li,b}).
\]

Li--ligand basin:
\[
A_{\rm LiL}=\{q:n_{\rm Li,L}\ge n_L^{\rm cut}\}.
\]

Coordination free energy:
\[
U_{\rm LiL}(q)
=
N_A\Delta g_{\rm LiL}^{\rm coord}
n_{\rm Li,L}(q).
\]

\subsection{Ligand competition with anion}

Li--anion coordination number:
\[
n_{\rm Li,A}(q)=\sum_{a\in A}s_{\rm Li,A}(r_{\rm Li,a}).
\]

Competition energy:
\[
U_{\rm comp}(q)
=
\frac12 k_{\rm comp}
\left[
n_{\rm Li,L}(q)+n_{\rm Li,A}(q)-n_{\rm shell}^{\rm target}
\right]^2.
\]

This shifts:
\[
c_{\rm free},
\quad
c_{\rm SSIP},
\quad
c_{\rm CIP},
\quad
c_{\rm LiL},
\quad
c_{\rm LiL\cdots A}.
\]

\subsection{Li--ligand--anion ternary state}

Define:
\[
A_{\rm LiL\cdots A}
=
\{
q:
n_{\rm Li,L}\ge n_L^{\rm cut},
\,
r_{\rm CIP}\le r_{\rm LiA}<r_{\rm SSIP}
\}.
\]

This state is distinct from:
\[
A_{\rm LiL}+A_{\rm free\ anion}.
\]

Its population is
\[
c_{\rm LiL\cdots A}
=
\int_{A_{\rm LiL\cdots A}}\rho(dq).
\]

\subsection{Additive-separated SSIP}

Define insertion coordinate
\[
\chi_{L{\rm sep}}(q)
=
\mathbf 1\{
L\text{ lies in the Li--anion corridor}
\}.
\]

The additive-separated SSIP basin is
\[
A_{\rm addSSIP}
=
A_{\rm SSIP}\cap\{\chi_{L{\rm sep}}=1\}.
\]

\subsection{Ligand-induced Li--anion covariance}

For state \(\alpha=A_{\rm addSSIP}\) or \(A_{\rm LiL\cdots A}\), define
\[
D_{\rm Li,A|\alpha}
=
\mathbb E[
v_{\rm Li}\cdot v_A
\mid A_\alpha
].
\]

The state charge mobility is
\[
D_Q(\alpha)
=
D_{\rm Li|\alpha}
+
D_{A|\alpha}
-
2D_{\rm Li,A|\alpha}.
\]

Ligand insertion may make
\[
D_{\rm Li,A|\alpha}<0,
\]
which increases \(D_Q\) and therefore conductivity.

\subsection{Additive microviscosity}

Local viscosity factor:
\[
f_L^\eta(q)
=
1+a_L^\eta n_{\rm Li,L}(q)+b_L^\eta \phi_L.
\]

Then
\[
\eta_{\rm loc}(q)=\eta_{\rm base}(q)f_L^\eta(q).
\]

This modifies:
\[
D_x(q),
\quad
K_{ij},
\quad
A_{\alpha\gamma}.
\]

\subsection{Ligand-shell residence memory}

Define memory coordinate:
\[
\psi_L(q)=n_{\rm Li,L}(q)-\mathbb E[n_{\rm Li,L}\mid A_{s(q)}].
\]

It enters:
\[
A_{\alpha\gamma}
=
\int
\nabla\psi_\alpha^\top D\nabla\psi_\gamma\,\rho(dq),
\]
\[
h_{\alpha a}
=
\int
\nabla\psi_\alpha^\top D\nabla P_a\,\rho(dq).
\]

\subsection{Ligand bridge}

For two Li centers \(\ell_1,\ell_2\):
\[
A_{\rm bridge}
=
\{q:n_{\ell_1,L}\ge n_L^{\rm cut},\ n_{\ell_2,L}\ge n_L^{\rm cut}\}.
\]

Bridge states contribute through:
\[
c_i,\quad K_{ij},\quad D_i^{\rm self},\quad A,h.
\]
```

This resolves the FEC/VC issue in the only first-principles-compatible way: neutral additives are ligand feature classes; the required FEC/bulky-anion behavior is the conditional Li–anion covariance (D_{\rm Li,A|\alpha}) in ligand-conditioned SSIP/CIP/additive-separated states, not a molecule-name branch. The uploaded checklist explicitly says additive-separated SSIP must be representable, species names may only be registry lookup labels, and the Li–anion covariance must change under ligand shell + bulky/delocalized anion conditions. 

